\chapter{Dataset e preprocessing}

\section{Descrizione del dataset}

Il dataset utilizzato per il progetto contiene informazioni relative a rotte aeree internazionali e risulta composto, nella sua forma iniziale, da 67663 record e 9 attributi. Le variabili originarie descrivono principalmente la compagnia aerea, l’aeroporto di origine, l’aeroporto di destinazione e alcuni campi accessori relativi agli identificativi dei soggetti coinvolti, alla presenza di \textit{codeshare}, al numero di scali e al tipo di equipaggiamento impiegato.

Dal punto di vista della Social Network Analysis, non tutte le variabili disponibili risultano ugualmente rilevanti. Poiché l’obiettivo del progetto è costruire una rete di collegamenti tra aeroporti, l’attenzione è stata rivolta soprattutto alle informazioni necessarie a definire i nodi e le relazioni del grafo. In particolare, le variabili di interesse principale sono risultate essere:
\begin{itemize}
    \item la compagnia aerea (\texttt{airline});
    \item l’aeroporto di origine (\texttt{source airport});
    \item l’aeroporto di destinazione (\texttt{destination airport}).
\end{itemize}

Le altre colonne presenti nel dataset sono state considerate utili soprattutto in fase di controllo e pulizia preliminare, ma non indispensabili per la costruzione della tabella finale utilizzata nelle successive analisi di rete.

\section{Pulizia e normalizzazione dei dati}

La prima operazione effettuata ha riguardato la normalizzazione dei nomi delle colonne. Nel dataset originale erano infatti presenti spazi superflui e una denominazione non uniforme di alcuni attributi. Per migliorare la leggibilità del codice e facilitare le elaborazioni successive, le colonne sono state rinominate con etichette più coerenti e standardizzate. In particolare, sono stati introdotti nomi come \texttt{airline\_id}, \texttt{source\_airport}, \texttt{source\_airport\_id}, \texttt{destination\_airport} e \texttt{destination\_airport\_id}.

Successivamente è stato eseguito un controllo dei valori mancanti. L’analisi ha mostrato che le colonne fondamentali per il progetto non presentano valori nulli: \texttt{airline}, \texttt{source\_airport} e \texttt{destination\_airport} risultano infatti sempre valorizzate. I campi con valori mancanti riguardano prevalentemente variabili accessorie, come \texttt{codeshare}, che presenta 53066 valori mancanti, ed \texttt{equipment}, che ne presenta 18. Poiché tali attributi non sono centrali rispetto agli obiettivi del progetto, la loro incompletezza non compromette la costruzione della rete.

Un ulteriore controllo ha riguardato la validità degli identificativi associati a compagnie e aeroporti. È stato osservato che in alcune righe gli ID non sono espressi come valori numerici validi, ma tramite il placeholder \texttt{\textbackslash N}, utilizzato nel dataset per rappresentare dati non disponibili o non correttamente specificati. In particolare:
\begin{itemize}
    \item la colonna \texttt{airline\_id} presenta 479 valori non validi;
    \item la colonna \texttt{source\_airport\_id} presenta 220 valori non validi;
    \item la colonna \texttt{destination\_airport\_id} presenta 221 valori non validi.
\end{itemize}

Poiché tali anomalie riguardano gli identificativi numerici e non i codici testuali degli aeroporti, non è stato necessario eliminare direttamente queste righe nella fase iniziale. Infatti, per la modellazione della rete risultano più importanti i codici degli aeroporti e della compagnia aerea, che sono risultati presenti e coerenti.

\section{Controllo di coerenza delle variabili principali}

Dopo la pulizia iniziale, è stata effettuata una verifica di coerenza sulle variabili che avrebbero costituito i nodi e le relazioni della rete. I risultati mostrano che:
\begin{itemize}
    \item la variabile \texttt{airline} presenta 568 valori distinti;
    \item la variabile \texttt{source\_airport} presenta 3409 valori distinti;
    \item la variabile \texttt{destination\_airport} presenta 3418 valori distinti.
\end{itemize}

Inoltre, è stato osservato che i codici degli aeroporti sono sempre composti da tre caratteri, caratteristica coerente con la presenza di codici aeroportuali sintetici utilizzabili come etichette dei nodi nel grafo. Anche i codici delle compagnie aeree risultano nella quasi totalità di lunghezza pari a due caratteri, con un numero molto limitato di casi a tre caratteri. Nel complesso, il dataset si presenta quindi sufficientemente pulito e coerente per essere trasformato in una struttura relazionale adatta alla costruzione della rete.



\section{Costruzione della tabella finale}

Una volta completati i controlli preliminari, è stata costruita una tabella finale ridotta, composta esclusivamente dagli attributi necessari alle analisi successive:
\begin{itemize}
    \item \texttt{source};
    \item \texttt{destination};
    \item \texttt{airline}.
\end{itemize}

La Tabella~\ref{tab:preprocessing_summary} riassume i principali risultati emersi durante la fase di preprocessing, includendo la dimensione iniziale del dataset, i controlli sui valori mancanti, la presenza di identificativi non validi e le caratteristiche della tabella finale ottenuta.

\begin{table}[H]
\centering
\caption{Sintesi dei principali controlli effettuati nella fase di preprocessing del dataset.}
\label{tab:preprocessing_summary}
\begin{tabular}{lcc}
\toprule
\textbf{Controllo} & \textbf{Variabile} & \textbf{Risultato} \\
\midrule
Dimensione iniziale del dataset & Record $\times$ attributi & 67663 $\times$ 9 \\
Valori mancanti & \texttt{airline} & 0 \\
Valori mancanti & \texttt{source\_airport} & 0 \\
Valori mancanti & \texttt{destination\_airport} & 0 \\
Valori mancanti & \texttt{codeshare} & 53066 \\
Valori mancanti & \texttt{equipment} & 18 \\
ID non validi & \texttt{airline\_id} & 479 \\
ID non validi & \texttt{source\_airport\_id} & 220 \\
ID non validi & \texttt{destination\_airport\_id} & 221 \\
Valori distinti & \texttt{airline} & 568 \\
Valori distinti & \texttt{source\_airport} & 3409 \\
Valori distinti & \texttt{destination\_airport} & 3418 \\
Dimensione tabella finale & Record $\times$ attributi & 67663 $\times$ 3 \\
Duplicati esatti & \texttt{(source, destination, airline)} & 0 \\
\bottomrule
\end{tabular}
\end{table}

La tabella finale mantiene 67663 record, corrispondenti alle rotte disponibili nel dataset originario, ma riorganizzati in una forma più adatta alla modellazione del grafo. I valori mancanti nelle colonne principali risultano nulli, e non sono presenti duplicati esatti rispetto alla tripla \texttt{(source, destination, airline)}. Questo significa che ogni riga identifica in modo univoco una specifica rotta operata da una determinata compagnia aerea.

Il risultato di questa fase è stato salvato nel file \texttt{routes\_prepared.csv}, che costituisce la base dati utilizzata nelle fasi successive del progetto per la costruzione del grafo diretto e della sua versione non diretta derivata.